\section{Introduction}
Dans le cadre de ce laboratoire, nous avons étudié les propriétés des matériaux piézoélectriques, en particulier leur comportement sous l'effet de tensions électriques. Les matériaux piézoélectriques sont largement utilisés dans des applications nécessitant des déplacements précis comme l'alignement de fibres optiques ou les microscopes à force atomique. Leur comportement est caractérisé par des phénomènes tels que l'hystérésis et le fluage qui influencent leur réponse à une tension appliquée.
\\\\
Pour analyser ces propriétés, nous avons utilisé deux approches expérimentales. La première consiste à mesurer l'hystérésis et du fluage en utilisant un système optique basé sur un laser He-Ne, une lame flexible et une photodiode à quatre quadrants dont nous avons mesuré la réponse du piézoélectrique à des signaux triangulaires et carrés. Cela nous a permis de quantifier la non-linéarité et les constantes de temps associées au fluage. La seconde méthode nous a permis de mesurer des déplacements nanométriques avec une grande précision. En observant les franges d'interférence, nous avons pu déterminer le coefficient piézoélectrique $d_{33}$ du matériau étudié.
\\\\
Nous avons alors déterminé que le temps de montée de la courbe de fluage est légèrement plus élevée que le temps de descente, que le coefficient de non-linéarité augmente avec la tension et la fréquence et que l'interférométrie avec le montage de Michelson permet de mesurer des nanodéplacements jusqu'à 31,5 nm avec un pas de 63 nm/V.

\section{Hystérésis et fluage}

\subsection{Description du dispositif}
Notre montage repose sur un système optique et électronique conçu pour étudier le comportement d’un matériau piézoélectrique. Le laser He-Ne (de longueur d’onde 632,8 nm) émet un faisceau rouge (voir Fig. 1) qui est dirigé vers une lame flexible polie comme un miroir. Cette lame est fixée à un actuateur piézoélectrique capable de se déformer sous l’effet d’une tension électrique.
\\\\
Lorsqu’une tension est appliquée à l’actuateur, ce dernier se contracte ou s’allonge. Cette différence de l'allongement fait fléchir la lame flexible. Cette déformation modifie l’angle du miroir et par conséquent la direction du faisceau laser réfléchi. Le faisceau dévié est ensuite capté par une photodiode à quatre quadrant dont le rôle est de mesurer avec précision la déviation du laser en convertissant la position du spot lumineux en un signal électrique. Plus la lame se courbe, plus le faisceau est dévié, et plus le signal généré par la photodiode varie.
\\\\
Pour piloter l’actuateur, on utilise un générateur de fonctions qui peut produire des signaux triangulaires ou carrés. Afin d'étudier le coefficient de non-linéarité pour un échantillon piezzo-électrique, nous mesurons sa courbe d'hystérèse pour plusieurs tensions (1V,2V,3V) à différentes fréquences à chaque fois (5Hz,10Hz,20Hz). Nous veillons à appliquer un offset de 1,5 V pour éviter les tensions négatives, car elles peuvent dépolariser le matériau piézoélectrique. Une dépolarisation du matériaux altère ses propriétés et réduit sa durée de vie. 
\\\\
Un boîtier électronique amplifie les signaux générés pour atteindre les tensions nécessaires au fonctionnement du piézo. Enfin, une application de mesure (PicoScope) sur l'ordinateur relié aux signaux amplifiés permet de visualiser en temps réel le signal de consigne (tension envoyée par le générateur) et le signal de mesure (réponse du piézo, via la photodiode). En superposant ces deux signaux, on peut tracer la courbe déplacement en fonction de la tension $d(V)$ et observer des phénomènes comme l’hystérésis ou le fluage.

\begin{figure}[H]
    \centering
    \includegraphics[scale=0.4]{schema_hysterese.png}
    \label{fig:schema_hysterese}
    \caption{Schéma du dispositif de mesure de l'hystérésis.}
\end{figure}

\subsection{Rappels théoriques}

\subsubsection{Hystérésis}
L'hystérésis dans les matériaux piézoélectriques résulte de la réorientation des domaines ferroélectriques sous champ électrique, ainsi que des frottements internes et des inerties mécaniques. Ce phénomène se manifeste par une boucle dans la courbe $d(V)$, où la réponse du matériau dépend de l’histoire des tensions appliquées. Le coefficient de non-linéarité est calculé comme suit :
\begin{equation}
    \text{Non-linéarité (\%)} = \frac{\text{Écart max (C-D)}}{\text{Plage totale (A-B)}} \times 100
    \label{eq:eq_lin}
\end{equation}
où A-B représente le déplacement total entre les points extrêmes de la boucle, et C-D représente l’écart maximal par rapport à une réponse linéaire idéale (voir Figure 1 du protocole).

\subsubsection{Fluage}
Le fluage décrit la réponse temporelle du piézoélectrique à un saut de tension. Il est modélisé par des équations exponentielles pour la montée et la descente :
\begin{equation}
    y(t) = a \left(1 - e^{-t/\tau}\right) + b \quad \text{(montée)}
    \label{eq:t_m}
\end{equation}
\begin{equation}
    y(t) = a e^{-t/\tau} + b \quad \text{(descente)}
    \label{eq:t_d}
\end{equation}
où $\tau$ est la constante de temps, typiquement de l’ordre de la milliseconde à la seconde.

\subsection{Résultats et Analyse}

\subsubsection{Hystérésis}
Nous pouvons observer les courbes d’hystérésis obtenues pour des tensions de 1 V, 2 V et 3 V (voir Figs. \ref{fig:1V_hyst}, \ref{fig:2V_lin} et \ref{fig:3V_lin}). Ces dernières montrent une boucle caractéristique dont la largeur augmente avec la tension appliquée (voir aussi sur Tab. \ref{tab:hysteresis_results_sorted}). Cette observation confirme que la non-linéarité du matériau piézoélectrique est plus marquée à haute tension.

\begin{figure}[H]
    \centering
    \includegraphics[scale=0.5]{1_1V.png}
    \caption{Courbe d'hystérésis pour une tension de 1 V.}
    \label{fig:1V_hyst}
\end{figure}

\begin{figure}[H]
    \centering
    \includegraphics[scale=0.5]{2_2V.png}
    \caption{Courbe $d(V)$ pour une tension de 2 V.}
    \label{fig:2V_lin}
\end{figure}

\begin{figure}[H]
    \centering
    \includegraphics[scale=0.5]{3_3V.png}
    \caption{Courbe $d(V)$ pour une tension de 3 V.}
    \label{fig:3V_lin}
\end{figure}

Les résultats du Tableau \ref{tab:hysteresis_results_sorted} montrent que la non-linéarité augmente avec la fréquence et la tension. Par exemple, pour une tension de 1 V, la non-linéarité passe de 0,328 \% à 5 Hz à 0,760 \% à 20 Hz. Cette tendance s’accentue pour des tensions plus élevées : à 3 V et 20 Hz, la non-linéarité atteint 0,765 \%. Cela s’explique par le fait que les domaines ferroélectriques du matériau ont moins de temps pour s’aligner à haute fréquence. Cela est ce qui augmente les frottements internes et donc la non-linéarité définit précédement par la relation Eq. \ref{eq:eq_lin}.

\begin{table}[h]
\centering
\begin{tabular}{|c|c|c|c|c|}
\hline
    \textbf{U / V} & \textbf{f / Hz} & \textbf{AB $\pm$ 0.01 V} & \textbf{CD $\pm$ 0.06 V} & \textbf{Non-linéarité $\pm$ 0.027 \%} \\
\hline
    1 & 5 & 0.022 & 0.007 & 0.328 \\
\hline
    2 & 5 & 0.046 & 0.017 & 0.381 \\
\hline
    3 & 5 & 0.056 & 0.024 & 0.421 \\
\hline
    1 & 10 & 0.020 & 0.011 & 0.519 \\
\hline
    2 & 10 & 0.036 & 0.021 & 0.571 \\
\hline
    3 & 10 & 0.050 & 0.029 & 0.587 \\
\hline
    1 & 20 & 0.015 & 0.012 & 0.760 \\
\hline
    3 & 20 & 0.039 & 0.030 & 0.765 \\
\hline
    2 & 20 & 0.028 & 0.022 & 0.783 \\
\hline
\end{tabular}
    \caption{Résultats des mesures d'hystérésis pour différentes tensions et fréquences, classés par ordre croissant de non-linéarité. Les incertitudes sur les mesures des tensions $U$ et des fréquences $f$ sont négligeables.}
    \label{tab:hysteresis_results_sorted}
\end{table}

\end{figure}

\newpage
\subsubsection{Fluage}
En appliquant un signal carré au système, nous mesurons la réponse temporelle du piézoélectrique (voir Fig. \ref{fig:fluage}). Les courbes de montée et de descente ont été ajustées avec les équations Eqs. \ref{eq:t_m} et \ref{eq:t_d}.

\begin{figure}[H]
    \centering
    \includegraphics[scale=0.45]{fluage_fit_montée_descente.png}
    \caption{Courbe de fluage montrant la réponse temporelle du piézoélectrique à un saut de tension.}
    \label{fig:fluage}
\end{figure}

On constate que l'ajustement correspond fortement au mesures. Les résultats de l’ajustement exponentiel sont présentés dans le Tableau \ref{tab:fluage_res} et nous permet d'affirmer que la constante de temps $\tau$ pour la montée, avec $\tau_m \approx$ 10,5 ms, est légèrement inférieure à celle de la descente avec $\tau_d \approx$ 10,8 ms (voir Tab. \ref{tab:fluage_res}). Cela signifie que le piézoélectrique atteint plus rapidement sa position finale lors de la montée que lors de la descente.

\begin{table}[h]
\centering
\begin{tabular}{|c|c|c|}
\hline
\textbf{Paramètre} & \textbf{Montée} & \textbf{Descente} \\
\hline
$A$ & $8.651228 \times 10^{-2} \pm 4.91 \times 10^{-5}$ & $8.352762 \times 10^{-2} \pm 3.15 \times 10^{-5}$ \\
\hline
$\tau$ (s) & $1.051606 \times 10^{-2} \pm 1.44 \times 10^{-5}$ & $1.078764 \times 10^{-2} \pm 1.00 \times 10^{-5}$ \\
\hline
$B$ & $-3.102036 \times 10^{-2} \pm 5.33 \times 10^{-5}$ & $-1.933693 \times 10^{-2} \pm 1.82 \times 10^{-5}$ \\
\hline
\end{tabular}
\caption{Résultats de l'ajustement exponentiel pour la montée et la descente.}
\label{tab:fluage_res}
\end{table}

Cette différence entre les constantes de temps de montée et de descente peut s’expliquer par les effets d’hystérésis internes. Lors de la montée, les domaines ferroélectriques s’alignent rapidement sous l’effet du champ électrique alors que lors de la descente les frottements internes et les inerties mécaniques ralentissent le retour à l’état initial. Cela montre que le comportement du piézoélectrique n’est pas parfaitement symétrique entre la charge et la décharge.
